\section{Step 21 Split a signed chroma displacement}\label{h263-prediction-s01}
Work through both components of $(-1,-1)$ before writing a prediction loop.
Divide the chroma displacement into an integer origin and a half-sample phase.
Keep the signed source value in a full-width register. Arithmetic shift by one
rounds a negative odd value down; masking its low bit yields phase one.

Apply the origin shift to the chroma macroblock coordinate, then form the byte
address with the chroma stride. The phase does not change the integer address
again. The public preparation helper supplies these already positioned windows;
you will retain this calculation in the eventual macroblock caller.

\paragraph{Observe} Luma component $-1$ gives origin shift $-1$, phase one.
Component $+1$ gives origin shift zero, phase one. Both use interpolation, but
start from different integer samples.
\Needspace{8\baselineskip}
\paragraph{Check} Which decomposition represents a chroma shift of $-0.5$?
\begin{itemize}
\item[A.] Origin zero and phase one.
\item[B.] Origin $-1$ and phase one.
\item[C.] Origin $-1$ and phase zero.
\end{itemize}

\section{Step 22 Admit the interpolation footprint}\label{h263-prediction-s02}
Validate the request and result records, then the fields and image extents.
Determine the reference width and height from the phase bits: either eight or
nine in each dimension. Check the reference stride against that width. Stride
eight is valid for phase zero horizontally, but cannot hold nine distinct active
samples in one row when the horizontal phase is one.

A halo is input data supplied by the caller. It is not permission to read all
bytes between active rows. Keep result bytes unchanged on every argument
rejection. Returning invalid status after a speculative vector read is too late.

\paragraph{Observe} Phases $(1,1)$ require nine reference samples in each of
nine rows. The current footprint remains eight by eight.
\Needspace{8\baselineskip}
\paragraph{Check} What is the minimal reference extent for diagonal prediction?
\begin{itemize}
\item[A.] $8s_R+9$ bytes, spanning nine active rows.
\item[B.] Eight bytes in total.
\item[C.] The size of the full output record.
\end{itemize}

\section{Step 23 Copy the integer prediction first}\label{h263-prediction-s03}
Begin scalar mode with phases zero. For each active position, load the current
and reference samples as unsigned bytes. Copy the reference value to the predicted
byte area, subtract it from the current value and store the difference as a
halfword. Advance each input row by its own stride, while using tight output rows.

Equal inputs should produce predicted samples equal to the current samples and
sixty-four zero residuals. Current zero against reference 255 should produce
sixty-four residuals of $-255$, not wrapped byte values of one.

\paragraph{Observe} The negative maximum-difference case returns SAD 16\,320.
Its residual halfwords remain negative when inspected as signed values.
\Needspace{8\baselineskip}
\paragraph{Check} Which value belongs in a residual when current is zero and prediction is 255?
\begin{itemize}
\item[A.] Positive 255.
\item[B.] Unsigned byte one.
\item[C.] Signed halfword $-255$.
\end{itemize}

\section{Step 24 Add one fractional axis}\label{h263-prediction-s04}
For horizontal phase one, load $A$ and $B$, add them in a width that holds 510,
add one and shift right by one. For vertical phase one, use $A$ and $C$ instead.
Only one neighbour direction is active in each of these cases. Do not load the
unused lower or right neighbour merely because the memory happens to be mapped.

Use an odd sum to check rounding. Samples 100 and 101 predict 101. Truncating
the average predicts 100 instead, changing both the stored prediction and the
residual. These errors can accumulate in the later reconstruction loop.

\paragraph{Observe} The worked caller with phases $(1,0)$ or $(0,1)$ prints
\texttt{128 0 112}. Its first sample is unchanged, but the block has nonzero
residuals elsewhere.
\Needspace{8\baselineskip}
\paragraph{Check} What is the baseline average of 100 and 101 at a half-sample position?
\begin{itemize}
\item[A.] 100, always rounding down.
\item[B.] 101, after adding one before division by two.
\item[C.] 201, without division.
\end{itemize}

\section{Step 25 Round a diagonal prediction once}\label{h263-prediction-s05}
For two active phases, load all four neighbours and form their sum before adding
two and shifting by two. Retain the unrounded intermediate. Do not reuse a helper
that has already rounded a horizontal average unless it preserves the information
needed for the final exact result.

Check the four-neighbour counterexample from the figure, then a high-valued
pattern. A sum of four 255s is 1020 and still predicts 255 after the bias and shift.
Both examples distinguish the width of a sample from the width of an intermediate.

\paragraph{Observe} With three zeros and one one, diagonal prediction is zero.
Two rounds of upward pair averaging would incorrectly produce one.
\Needspace{8\baselineskip}
\paragraph{Check} When is the diagonal rounding bias applied?
\begin{itemize}
\item[A.] Once, after summing all four neighbours.
\item[B.] Separately after every addition.
\item[C.] Only after the signed residual has been clipped to a byte.
\end{itemize}

\section{Step 26 Gather eight samples without overreading}\label{h263-prediction-s06}
Implement an MSA row with eight halfword lanes. Load active bytes with unsigned
scalar loads and insert their values into the corresponding lanes. Repeat for
each neighbour vector that the chosen phases require. Every lane must be written
before an operation reads that vector; inactive neighbour vectors need not be
initialized because their branches must not consume them.

This arrangement uses vector arithmetic without pretending that an eight-sample
row occupies sixteen input bytes. A different implementation may coalesce tight
rows or stage a bounded tile. Verify the complete access footprint, including
its final row. Safe access is part of the result contract, not an optional cost
of making the arithmetic faster.

\paragraph{Observe} For a padded row, the implementation reads only eight
current samples and the selected eight or nine reference columns. A tight-row
coalesced load is valid only when all of its bytes are active samples.
\Needspace{8\baselineskip}
\paragraph{Check} Why can a sixteen-byte load be wrong for this current row?
\begin{itemize}
\item[A.] MSA cannot calculate unsigned values.
\item[B.] All byte pointers must be sixteen-byte aligned.
\item[C.] The load can include padding or samples outside the admitted footprint.
\end{itemize}

\section{Step 27 Keep interpolation and residual lanes wide}\label{h263-prediction-s07}
Apply the phase-specific arithmetic to the halfword neighbour vectors. A rounded
unsigned average is sufficient for one fractional axis. The diagonal path sums
four inputs and performs one rounded shift. Subtract the predicted halfwords
from the current halfwords, preserving signed difference bit patterns.

Store the eight residual halfwords as one complete sixteen-byte output row.
Extract and store the eight predicted bytes without overwriting neighbouring
output data. Writing sixteen prediction bytes for an eight-byte row can corrupt
a residual row that was already produced, even if the extra write remains inside
the overall record.

\paragraph{Observe} Byte orders change the two memory bytes of a signed
halfword, but not its decoded value or the predicted sample sequence.
\Needspace{8\baselineskip}
\paragraph{Check} Which width holds a diagonal neighbour sum before rounding?
\begin{itemize}
\item[A.] An unsigned halfword, whose range includes 1020.
\item[B.] A byte, because each input fits a byte.
\item[C.] A one-bit phase field.
\end{itemize}

\section{Step 28 Accumulate a score without erasing signs}\label{h263-prediction-s08}
Form a separate magnitude vector from the signed residual vector. Accumulate
those magnitudes in halfwords across the eight rows, then widen and reduce to
the scalar return value. Each lane accumulates at most $8\times255=2040$;
the full block reaches 16\,320. Keep the stored signed residuals unchanged.

Check $P+e=C$ independently at every output position. This catches reversed
subtraction even when the absolute sum still agrees. It also catches a misplaced
halfword or an incorrect native-endian store that a total alone could hide.

\paragraph{Observe} Positive and negative maximum-difference blocks have the
same SAD, but their residual arrays have opposite signs.
\Needspace{8\baselineskip}
\paragraph{Check} Why does a correct SAD not prove a correct residual array?
\begin{itemize}
\item[A.] SAD specifies the required motion-vector syntax.
\item[B.] Taking magnitudes removes the direction of every difference.
\item[C.] Signed halfwords cannot be read in both byte orders.
\end{itemize}

\section{Step 29 Return a complete block to the caller}\label{h263-prediction-s09}
Publish success after all sixty-four predictions and residuals are stored.
Return the magnitude sum, restore the outer caller's saved state and stop at its
continuation in Debug. Inspect the first predicted byte, the first signed
halfword at result offset 64 and the two return registers before continuing.

Prepare the four phase combinations as separate new files. The default diagonal
example prints \texttt{120 8 232}; the integer example prints \texttt{128 0 0}.
They share origins but choose different interpolation positions. The samples
come from raw Cb planes for this arithmetic observation; the completed encoder
must supply its locally reconstructed reference.

\paragraph{Observe} Development status four prints the unfinished message and
exits one. Completed scalar and MSA modes produce the same records and successful
exit status for each worked phase pair.
\Needspace{8\baselineskip}
\paragraph{Check} At the caller return boundary, what is already complete?
\begin{itemize}
\item[A.] A conforming whole-video output file.
\item[B.] The future quantized coefficient stream.
\item[C.] This block's prediction/residual record and returned status/SAD.
\end{itemize}

\section{Step 30 Check every sample before measuring work}\label{h263-prediction-s10}
Run completed feedback against your current source:
\begin{verbatim}
python3 construction/h263-prediction-construction/grade.py \
    my-prediction.asm --cli /path/to/mars-cli.jar \
    --report prediction-results.json
\end{verbatim}
The 96 public cases cover all phases, both modes and byte orders, positive and
negative extremes, equal samples, mixed patterns, a diagonal rounding counterexample
and invalid requests. The checker calculates every expected sample independently,
checks signed residuals and the exact reconstruction identity, and verifies input,
request, output-guard and saved-state contracts. Reports require new paths.

Measure from entry through return, including admission, gathers, arithmetic,
stores and reduction. Counting vector arithmetic alone omits the cost of placing
samples into lanes. Preserve the same footprint and exact rounding when optimizing.
The returned score and instruction count describe this block, not a completed
encoder's quality or speed.

\paragraph{Observe} A safe alternative can use different rounding instructions
or a larger private frame. It must retain the same sample results and bounded
accesses. Later transform work will consume the signed halfwords produced here.
\Needspace{8\baselineskip}
\paragraph{Check} What must a faster gather preserve?
\begin{itemize}
\item[A.] The reference solution's exact scratch-register names.
\item[B.] Every required sample and the admitted access footprint.
\item[C.] Only the final absolute sum, regardless of signed residuals.
\end{itemize}
